\documentclass[11pt]{article}
%autocompile publish
\usepackage{math-hse,array,tikz,amsfonts}
\usetikzlibrary{hobby,decorations.pathreplacing,arrows.meta}
\title{Лекция 2. Основные понятия теории вероятностей}

\date{??.01.2027}
\begin{document}
\section{Формула полной вероятности}
Допустим, в некотором вузе есть три факультета: математический, физический и
химический. Факультеты отличаются по числу студентов~--- например, на математическом
факультете учится $1/10$ студентов всего вуза, на физическом~--- $3/10$, а на
химическом~--- оставшиеся $6/10$ (то есть $3/5$) всех студентов вуза (см.
табл.~\ref{otlichniki}).\footnote{Конечно, мы считаем, что каждый студент учится
только на одном факультете, и каждый учится на каком-нибудь факультете из этих
трёх.} Мы также знаем, какова доля отличников на каждом факультете. Например, на
математическом факультете каждый четвёртый является отличником, на физическом~---
каждый второй студент отличник, а на химическом $3/4$ студентов являются
отличниками. Какова вероятность, что случайно взятый студент вуза является
отличником?

\begin{table}[hbt]
	\begin{center}
		\begin{tabular}{|c|c|c|c|c|c|c|}
			\hline
			факультет & матфак & физфак & химфак \\
			\hline
			какая часть студентов учится на факультете & $1/10$  & $3/10$ & $6/10$ \\
			\hline
			доля отличников на факультете & $1/4$ & $1/2$ & $3/4$\\
			\hline
		\end{tabular}
		\caption{Распределение студентов по факультетам}\label{otlichniki}
	\end{center}
\end{table}

Допустим, что всего в вузе учится $N$ студентов. Тогда на математическом факультете
будет $N\cdot \frac{1}{10}$ студентов. Среди них четвертая часть (то есть
$N\cdot \frac{1}{10} \cdot \frac{1}{4}=N\cdot \frac{1}{40}$ человек) являются
отличниками. То есть вероятность, что случайно выбранный студент является отличником
с математического факультета, равна
$N \cdot \frac{1}{10} \cdot \frac{1}{4} \cdot \frac{1}{N}
= \frac{1}{10} \cdot \frac{1}{4}=0,025$.

К тому же выводу можно прийти и через теорию вероятностей. Обозначим событие, что
случайно взятый студент~--- математик, за $H_1$, а событие, что случайно взятый
студент~--- отличник, за $A$. В условии дано $P(H_1) = \frac{1}{10}$. Выберем
случайного студента среди студентов математического факультета. В таком случае
вероятность, что он окажется отличником~--- это \emph{условная вероятность события
$A$ при условии $H_1$} (то есть $P(A|H_1)$). Эта вероятность в точности равна доле
отличников среди студентов математического факультета (то есть $P(A \vert H_1)
= \frac{1}{4}$). В то же время, по определению условной вероятности,
\begin{equation}
	P(A|H_1)=\frac{P(A\cap H_1)}{P(H_1)},
\end{equation}
где $P(A\cap H_1)$~--- вероятность одновременного выполнения $A$ и $H_1$. Умножая
обе части этого равенства на $P(H_1)$, получим:
\begin{equation}\label{eq:AH1}
	P(A\cap H_1)=P(A|H_1)\cdot P(H_1).
\end{equation}
Подставляем известные числа и получаем тот же результат: $P(A \cap H_1) =
\frac{1}{10} \cdot \frac{1}{4} = \frac{1}{40} = 0,025$.

Обозначим событие, что случайно взятый студент~--- физик, за $H_2$, а событие, что
случайно взятый студент~--- химик, за $H_3$. Тогда вероятности, что случайный
студент окажется физиком-отличником либо химиком-отличником, вычисляются аналогично
формуле~(\ref{eq:AH1}):
$$
    P(A \cap H_2) = P(H_2) \cdot P(A \vert H_2) = \frac{3}{10} \cdot \frac{1}{2}
= \frac{3}{20} = 0,15
$$
$$
    P(A \cap H_3) = P(H_3) \cdot P(A \vert H_3) = \frac{6}{10} \cdot \frac{3}{4}
= \frac{18}{40} = 0,45
$$

Один студент учится на одном и только на одном факультете, то есть события
$H_1$, $H_2$, $H_3$ попарно несовместны и вместе образуют всё вероятностное
пространство. Иными словами, эти события разбивают всё вероятностное пространство
на несколько непересекающихся частей. В этом случае такие события часто называют
\emph{гипотезами}, и говорят, что они образуют <<полную систему событий>>. Тот факт,
что гипотезы попарно несовместны, и что в объединении они дают всё вероятностное
пространство, очень важен~--- он означает, что всегда реализуется в точности одна
гипотеза.

Событие $A$ можно представить как объединение кусочков, лежащих в $H_1$, $H_2$,
$H_3$, то есть $A=(A\cap H_1)\cup (A\cap H_2)\cup (A\cap H_3)$. События
$A\cap H_1$, $A\cap H_2$, $A\cap H_3$ тоже не пересекаются, и значит их вероятности
можно складывать: $P(A)=P((A\cap H_1)\cup (A\cap H_2)\cup (A\cap H_3))
=P(A\cap H_1)+P(A\cap H_2)+P(A\cap H_3)$. Учитывая~(\ref{eq:AH1}), имеем:
\begin{equation}\label{eq:fp}
    P(A)=P(A|H_1)\cdot P(H_1)+P(A|H_2)\cdot P(H_2)+P(A|H_3)\cdot P(H_3)
\end{equation}

Доказанная нами формула~(\ref{eq:fp}) называется \textit{формулой полной
вероятности}. В общем случае гипотез может быть не три, а больше.

\paragraph{Формула полной вероятности} Если имеется:
\begin{enumerate}
    \item Вероятностное пространство $\Omega$
    \item Полная системы событий (гипотез) $H_{1}, H_{2},..., H_{n} \subset \Omega$,
    т.~е. такие гипотезы, что $$H_1 \sqcup H_2 \sqcup \dots \sqcup H_n = \Omega$$
    \item Произвольное событие $A \subset \Omega$
\end{enumerate}
то вероятность события $A$ (полная вероятность) вычисляется по формуле:
$$P(A)=P(H_{1})P(A|H_{1})+P(H_{2})P(A|H_{2})+\dots+P(H_{n})P(A|H_{n})$$

\begin{example}

    На столе лежат три монетки: две обыкновенные и одна с орлами с двух сторон.
    Наудачу выбирается одна монетка и подбрасывается. С какой вероятностью выпадет
    орёл?

    Мы могли выбрать либо обычную монетку, либо монетку с двумя орлами. В каждом
    из этих случаев несложно найти вероятность выпадения орла, поэтому обозначим
    за $H_{1}$ событие (гипотезу) <<мы взяли обычную монетку>>, за $H_{2}$~---
    событие (гипотезу) <<мы взяли монетку с двумя орлами>>. Интересующее нас
    событие $A$~--- выпадение орла на выбранной монетке. Итак, имеем:

\vspace{1em}

$P(H_{1})=\frac{2}{3}$,~~~~$P(A|H_{1})=\frac{1}{2}$

\vspace{1em}

$P(H_{2})=\frac{1}{3}$,~~~~$P(A|H_{2})=1$

\vspace{1em}

Значит, $$P(A)=P(H_{1})P(A|H_{1})+P(H_{2})P(A|H_{2})=\frac{2}{3}\cdot\frac{1}{2}+\frac{1}{3}\cdot1=\frac{2}{3}.$$

\end{example}

\section{Формула Байеса}

Вернемся к задаче об отличниках. Пусть мы выбрали случайным образом студента, и он
оказался отличником. Теперь нас интересует вероятность того, что он учится на
математическом факультете. Итак, задача в отыскании условной вероятности того, что
студент учится на первом факультете, при условии того, что он отличник. Известны
вероятности $P(H_{1})$,~$P(H_{2})$,~$P(H_{3})$ того, что студент учится на каждом
из трёх факультетов, и вероятности (условные!) $P(A|H_{1}),~P(A|H_{2}),~P(A|H_{3})$
того, что студент, учась на каждом из этих факультетов, будет отличником. Нас
интересует условная вероятность $P(H_{1}|A)$, что студент, \textit{оказавшийся}
отличником, учится на математическом факультете.

Напомним, что по формуле полной вероятности, $P(A)$ представляется в виде суммы
трех слагаемых, каждое из которых выражает долю отличников с соответствующего
факультета от числа всех студентов вуза:
\begin{equation*}
	P(A)=P(A|H_1)\cdot P(H_1)+P(A|H_2)\cdot P(H_2)+P(A|H_3)\cdot P(H_3)
\end{equation*}
% В этом месте есть кажущееся несоответствие условию задачи. Ведь мы \textit{знаем},
% что студент оказался отличником, т.\,е.~ событие $A$ произошло, зачем же мы ищем
% вероятность этого события? На самом деле, здесь $P(A)$ \textit{априорная}
% вероятность события $A$, просто вероятность выбрать отличника, выбирая случайного
% студента.

Нас интересует, какую долю составляет первое слагаемое ($P(A|H_1)\cdot P(H_1)$) в
этой сумме~--- это и есть условная вероятность выбрать отличника с математического
факультета, выбирая из всех отличников. Иными словами, нас интересует отношение
\begin{equation*}
	P(H_1|A)
    = \frac{P(A|H_1)\cdot P(H_1)}{P(A)}
    = \frac{P(A|H_1)\cdot P(H_1)}{P(A|H_1)\cdot
	P(H_1)+P(A|H_2)\cdot P(H_2)+P(A|H_3)\cdot P(H_3)}
\end{equation*}
Полученная формула называется \emph{формулой апостериорной вероятности}, или
\emph{формулой Байеса} В рассмотренном примере интересующая нас вероятность есть
отношение $\frac{0.025}{0.625}=0.04$. Как и в случае формулы полной вероятности,
формула Байеса может применяться для произвольного числа гипотез.

\paragraph{Формула Байеса} Если имеется:
\begin{enumerate}
    \item Вероятностное пространство $\Omega$
    \item Полная системы событий (гипотез) $H_{1}, H_{2},..., H_{n} \subset \Omega$,
    т.~е. такие гипотезы, что $$H_1 \sqcup H_2 \sqcup \dots \sqcup H_n = \Omega$$
    \item Произвольное событие $A \subset \Omega$
\end{enumerate}
то вероятность гипотезы $H_i$ при условии события $A$ (апостериорная вероятность)
вычисляется по формуле:
$$
    P(H_i \vert A)
    = \frac{P(H_i)P(A|H_i)}
    {P(H_{1})P(A|H_{1})+P(H_{2})P(A|H_{2})+\dots+P(H_{n})P(A|H_{n})}
$$

\begin{examples}
\item[]
\begin{enumerate}
    \item На столе лежат три монетки: две обыкновенные и одна с орлами с двух сторон.
    Наудачу выбирается одна монетка и подбрасывается. На подброшенной монетке
    выпал орёл. С какой вероятностью мы кидали обычную монетку?
    
    Пусть:
    
    $H_{1}$ событие (гипотеза) ~--- мы взяли обычную монетку,
    
    $H_{2}$ событие (гипотеза) ~--- мы взяли монетку с двумя орлами.
    
    Гипотезы $H_{1}$ и $H_{2}$ образуют полную систему событий.
    
    Событие $A$~--- на монетке выпал орёл.
    
    Тогда
	
    \vspace{1em}

    $P(H_{1})=\frac{2}{3}$, ~~~~ $P(A|H_{1})=\frac{1}{2}$

    \vspace{1em}

    $P(H_{2})=\frac{1}{3}$, ~~~~ $P(A|H_{2})=1$

    \vspace{1em}

    Значит,
    $$
        P(H_1 \vert A)
        =\frac{P(H_{1})P(A|H_{1})}{P(H_{1})P(A|H_{1})+P(H_{2})P(A|H_{2})}
        =\frac{\frac{2}{3}\cdot\frac{1}{2}}
        {\frac{2}{3}\cdot\frac{1}{2}+\frac{1}{3}\cdot1}
        =\frac{1/3}{2/3}=\frac{1}{2}
    $$

    Это и неудивительно: хотя на странной монетке орлы встречаются в два раза чаще,
    чем на обычной, зато на нашем столе обычные монетки встречаются в два раза чаще,
    чем странные.

    \item Тест на ВИЧ выдаёт верный результат в $95\%$ случаев. Известно, что ВИЧ
    заражена $1/1000$ часть всего населения. Человек получил положительный
    результат теста на ВИЧ (то есть по результатам теста он заражён). С какой
    вероятностью он действительно заражён?
    
    Пусть
    
    $H_{1}$ событие (гипотеза) ~--- человек заражён ВИЧ,
    
    $H_{2}$ событие (гипотеза) ~--- человек не заражён ВИЧ.
    
    Гипотезы $H_{1}$ и $H_{2}$ образуют полную систему событий.
    
    Событие $A$ ~--- тест дал положительный результат.
    
    Тогда
    
    \vspace{1em}

    $P(H_{1})=0,001$,~~~~      $P(A|H_{1})=0,95 $
    
    $P(H_{2})=0,999$,~~~~      $P(A|H_{2})=0,05$ \footnote{
        В реальности $P(A \vert H_1)$ и $P(A \vert H_2)$ не обязательно будут
        давать в сумме $1$. Например, тест может верно определять заражённых людей
        в $95\%$ случаев, но при этом ошибочно определять здоровых людей как
        заражённых в $1\%$ случаев. В этом случае $P(A \vert H_1) = 0,95$, а
        $P(A \vert H_2) = 0,01$.
    }

    \vspace{1em}

    Значит,
    $$
        P(H_{1}|A)
        =\frac{P(H_{1})P(A|H_{1})}{P(H_{1})P(A|H_{1})+P(H_{2})P(A|H_{2})}
        =\frac{0,001\cdot0,95}{0,001\cdot0,95+0,999\cdot0,05}\approx0,0187
    $$
    
    Результат кажется удивительным: ведь ошибка теста всего $5\%$, как же могло
    получиться, что вероятность заболевания при положительном результате теста
    всего $1,87\%$? Дело в том, что среди людей, получивших положительный
    результат теста, гораздо больше получивших ложный положительный результат
    ($5\%$ от $99,9\%$ населения), чем людей, действительно зараженных ВИЧ
    ($0{,}1\%$ населения), а тем более, действительно больных, получивших
    положительный результат теста. % ~--- см. рисунок~\ref{fig:hivtest}.

    % \begin{figure}[htb]
    % 	\begin{center}
    % 		% \includegraphics[scale=0.7]{hivtest.pdf}
    % 		\caption{Применение формулы Байеса в случае теста на ВИЧ. Тонкая полоска слева~--- люди, зараженные ВИЧ. Темный закрашенный прямоугольник~--- те из них, кто получил положительный результат теста. Длинный закрашенный прямоугольник справа — здоровые люди, получившие положительный результат теста. Доля больных людей мала по сравнению со всеми людьми, получившими положительный результат.}\label{fig:hivtest}
    % 	\end{center}
    % \end{figure}

    \item Пусть имеется две коробки. В первой коробке $5$ белых и $4$ чёрных шара,
    во второй~--- $3$ белых и $2$ чёрных. Из второй коробки случайным образом
    выбирается один шар и перекладывается в первую коробку, после чего из первой
    коробки наудачу выбирается один шар. Какова вероятность того, что этот шар
    белый? Если оказалось, что этот шар белый, то с какой вероятностью из второй
    коробки переложили белый шар?

    Пусть:
    
    $H_{1}$ событие (гипотеза) ~--- переложили белый шар,
    
    $H_{2}$ событие (гипотеза) ~--- переложили чёрный шар.
    
    Гипотезы $H_{1}$ и $H_{2}$ образуют полную систему событий.
    
    Событие $A$~--- из первой коробки достали белый шар.
    
    Тогда
	
    \vspace{1em}

    $P(H_{1})=\frac{3}{5}$, ~~~~ $P(A|H_{1})=\frac{5+1}{9+1}=\frac{6}{10}$

    \vspace{1em}

    $P(H_{2})=\frac{2}{5}$, ~~~~ $P(A|H_{2})=\frac{5}{9+1}=\frac{5}{10}$

    \vspace{1em}

    Значит,
    $$
        P(A)=P(H_{1})P(A|H_{1})+P(H_{2})P(A|H_{2})
        =\frac{3}{5}\cdot\frac{6}{10}+\frac{2}{5}\cdot\frac{5}{10}
        =\frac{28}{50}=\frac{14}{25}
    $$
    $$
        P(H_1 \vert A)
        =\frac{P(H_{1})P(A|H_{1})}{P(A)}
        =\frac{\frac{3}{5}\cdot\frac{6}{10}}
        {\frac{14}{25}}
        =\frac{18/50}{28/50}=\frac{18}{28}=\frac{9}{14}
    $$

\end{enumerate}
\end{examples}

\end{document}
